\documentclass[11pt, a4paper]{article}

% --- Encodage et Langue ---
\usepackage[utf8]{inputenc}
\usepackage[T1]{fontenc}
\usepackage[french]{babel}

% --- Mise en page et Géométrie ---
\usepackage[margin=2.5cm]{geometry}
\usepackage{fancyhdr} % Pour les en-têtes et pieds de page
\usepackage{setspace}
\onehalfspacing % Interligne un et demi pour le confort de lecture

% --- Mathématiques ---
\usepackage{amsmath}
\usepackage{amssymb}
\usepackage{amsthm}
\usepackage{mathtools}

% --- Couleurs et Graphiques ---
\usepackage{xcolor}
\usepackage{graphicx}
\usepackage{tikz}

% --- Liens et Références ---
\usepackage{hyperref}
\hypersetup{
    colorlinks=true,
    linkcolor=blue!70!black,
    citecolor=blue!70!black,
    urlcolor=blue!70!black
}

% --- Style des en-têtes ---
pagestyle{fancy}
\fancyhf{}
\fancyhead[L]{\textit{Nom du cours / Module}}
\fancyhead[R]{\textit{Mes Notes - Date}}
\fancyfoot[C]{\thepage}

% --- Environnements de théorèmes (optionnel) ---
\theoremstyle{definition}
\newtheorem{defn}{Définition}[section]
\newtheorem{thm}{Théorème}[section]
\newtheorem{prop}{Proposition}[section]
\newtheorem{ex}{Exemple}[section]

% ==========================================
%           DEBUT DU DOCUMENT
% ==========================================
\begin{document}

% --- Page de garde / Titre ---
\begin{center}
    {\huge \textbf{Titre du Cours}} \\[0.5cm]
    {\large Sous-titre ou Chapitre} \\[0.5cm]
    \vspace{0.3cm}
    \textbf{Auteur :} Votre Nom \\
    \textbf{Date :} \today
\end{center}

\hrule
\vspace{1cm}

% --- Table des matières (optionnel) ---
\tableofcontents
\newpage

% ==========================================
section{Introduction}
% ==========================================
Rédigez votre introduction ici. C'est l'endroit idéal pour poser les bases du cours ou du chapitre.

% ==========================================
section{Concepts clés}
% ==========================================

\subsection{Première notion}
Vous pouvez insérer des formules mathématiques en ligne comme $E = mc^2$, ou bien en bloc centré :
$$
\int_{-\infty}^{+\infty} e^{-x^2} dx = \sqrt{\pi}
$$

\begin{defn}
    Une définition importante s'écrit de cette manière pour bien ressortir dans vos notes.
\end{defn}

% ==========================================
section{Conclusion}
% ==========================================
Résumé rapide des points à retenir pour les révisions.

\end{document}